\begin{abstract}
\noindent\textbf{Purpose} -- This paper asks where capacity in artificial intelligence (AI) for
climate adaptation should be built first when the welfare of poorer households counts for more,
and whether such capacity has so far made incomes less vulnerable to heat.

\smallskip\noindent\textbf{Design/methodology/approach} -- We develop equity-weighted doubly
robust policy learning (EW-DRPL). The method uses exogenous temperature shocks, orthogonal doubly
robust scores for how AI capability changes shock sensitivity, Atkinson welfare weights over income
deciles, and budget-constrained decision trees. We derive a regret bound, run simulations and
apply the method to cereal yields in 145 countries, measuring climate-AI capability with OpenAlex
data.

\smallskip\noindent\textbf{Findings} -- In simulations, EW-DRPL has less than one-third of the regret
of standard interaction regressions, and equity weighting reduces regret by 36--43 per cent under high
inequality aversion. Empirically, a 1$^{\circ}$C hotter-than-trend hottest month lowers
cereal-yield growth by 3.6 percentage points. Climate-AI capability has no detectable effect on this
damage, and effects above two-thirds of it are ruled out. Allocating capacity to existing AI leaders
yields almost no equity-weighted benefit.

\smallskip\noindent\textbf{Research limitations/implications} -- Bibliometric output measures
research capacity rather than deployment, and overlap between countries is limited.

\smallskip\noindent\textbf{Practical implications} -- Development partners should evaluate
AI-for-adaptation programmes against measured damages and prioritise heat-exposed, agriculturally
dependent settings over existing AI leaders.

\smallskip\noindent\textbf{Social implications} -- Making inequality aversion explicit shows how
targeting choices shift resilience gains towards the poor.

\smallskip\noindent\textbf{Originality/value} -- To our knowledge, the paper is the first to estimate equity-weighted
allocation rules for AI-enabled climate resilience, with open code and data.

\smallskip\noindent\textbf{Keywords} -- Artificial intelligence, Climate adaptation, Inequality,
Policy learning, Double machine learning, Agriculture

\smallskip\noindent\textbf{Paper type} -- Research paper
\end{abstract}
